\vspace{-1em}
\section{Conclusion}
\vspace{-1em}
In this paper, we evaluate multi-agent systems for their robustness and ability to successfully complete benign tasks. We introduce the TAMAS benchmark, which comprises 300 adversarial attack scenarios and 100 benign scenarios spanning five domains and six attack types. To understand how agent coordination affects vulnerability, we experiment with three agent interaction configurations. Our findings reveal that multi-agent frameworks are highly susceptible to adversarial attacks, highlighting the urgent need for stronger defense mechanisms to ensure their safety. Discussion of limitations and future work is provided in Appendix \ref{App:limitations}.

\section*{Reproducibility statement}
To ensure transparency and reproducibility, we are committed to making our research accessible. We provide comprehensive experimental details in the paper, and all datasets and code will be publicly released upon publication. All experiments were conducted using open-source frameworks AutoGen and CrewAI, with models accessed via API or Ollama.

\section*{Ethics statement}
This work investigates the robustness of multi-agent LLM systems against adversarial attacks. The primary aim of this study is to systematically evaluate how different system configurations and attack strategies influence both safety and task performance. Our findings are intended to advance the development of safer and more reliable multi-agent AI systems.





